% Part: normal-modal-logic
% Chapter: filtrations
% Section: S5-decidable

\documentclass[../../../include/open-logic-section]{subfiles}

\begin{document}

\olfileid[id]{nml}{fil}{dec}

\olsection{\Log{S5} Dapat Diputuskan}
Sifat model berhingga memberi kita cara mudah untuk menunjukkan bahwa sistem
logika modal yang diberikan oleh skema-skema dapat \emph{diputuskan} (yakni
bahwa terdapat prosedur yang dapat dihitung untuk menentukan apakah
!!a{formula} !!{derivable} dalam sistem tersebut atau tidak).

\begin{thm}
  \Log{S5} dapat diputuskan.
\end{thm}

\begin{proof}
  Misalkan $!A$ diberikan, dan andaikan !!{propositional variable}s yang
  muncul dalam $!A$ terdapat di antara $p_1$, \dots, $p_k$. Karena bagi
  setiap $n$ hanya terdapat berhingga banyak model universal dengan $n$ dunia
  yang menetapkan suatu nilai kepada $p_1$, \dots, $p_k$, kita dapat
  mengenumerasi, \emph{secara paralel}, semua teorema \Log{S5} dengan
  menghasilkan pembuktian secara sistematis, serta semua model universal dengan
  $1$, $2$, \dots dunia dan memeriksa apakah $!A$ gagal pada suatu
  dunia dalam salah satu model tersebut. Pada akhirnya, salah satu dari kedua
  proses paralel itu akan memberikan jawaban, sebab menurut
  \olref[com][fra]{thm:generaldet} dan \olref[fmp]{cor:S5fmp}, $!A$
  !!{derivable} atau gagal dalam suatu model universal berhingga.
\end{proof}

Pembuktian di atas berlaku bagi \Log{S5} karena filtrasi dari model universal
secara otomatis bersifat universal. Hal yang sama berlaku bagi refleksivitas
dan serialitas, tetapi sifat-sifat lain memerlukan lebih banyak upaya.

\end{document}
